% Readout 7: generate a rationale first, then read the answer-token distribution
% (contrast only; conflicts with the one-forward System One positioning).
% Prompt x:[B,T0]; sampled rationale r:[B,G] (G decode steps, KV cache);
% then readout 1 on [x; r; answer cue]. T = T0 + G + 1.
% MetisCraft readout tensor figures: shared notation, palette and macros.
% Cell counts are illustrative; each symbolic axis keeps one physical edge
% length inside a figure (u per cell): T=12, d=5, V=12, K=L=3, B=2, d_h=4.
\documentclass[tikz,border=3mm]{standalone}
\usepackage[UTF8,fontset=fandol]{ctex}
\usepackage{amsmath,amssymb,bm}
\usetikzlibrary{arrows.meta,calc,backgrounds,decorations.pathreplacing}
\definecolor{Ink}{HTML}{273648}
\definecolor{Muted}{HTML}{607086}
% token roles
\definecolor{Ctx}{HTML}{85898F}    % state / context
\definecolor{Qst}{HTML}{4F8FA5}    % question
\definecolor{Opt}{HTML}{5A9A6E}    % options / candidates (three lightness levels)
\definecolor{Gen}{HTML}{B89A5E}    % generated rationale tokens
% tensor roles
\definecolor{Hid}{HTML}{8A74B5}    % hidden activations
\definecolor{Wgt}{HTML}{E0904F}    % trainable / pretrained weights
\definecolor{Lgt}{HTML}{C95B5B}    % logits, probabilities, loss
\definecolor{Idx}{HTML}{7E8B9B}    % integer indices / labels
\definecolor{Msk}{HTML}{3F5A78}    % binary attention support
\colorlet{TokS}{Ctx!55}\colorlet{TokQ}{Qst!65}
\colorlet{Tok1}{Opt!30}\colorlet{Tok2}{Opt!58}\colorlet{Tok3}{Opt!88}
\colorlet{TokD}{Ink!80}\colorlet{TokG}{Gen!65}\colorlet{TokM}{Opt!100!black!60}
\colorlet{TokC}{Ink!80}\colorlet{TokP}{white}
\colorlet{TokL}{Opt!58}\colorlet{TokT}{Qst!65}
\newcommand{\U}{0.3}
\newcommand{\Off}{0.14}
\newcounter{tseed}
\tikzset{
  flow/.style={-{Latex[length=2mm,width=1.4mm]},draw=Ink!85,line width=.7pt},
  dflow/.style={-{Latex[length=2mm,width=1.4mm]},draw=Ink!60,line width=.6pt,dashed},
  op/.style={draw=Ink!30,fill=white,rounded corners=2pt,line width=.5pt,
             align=center,font=\small,inner sep=4pt},
  sym/.style={font=\small,inner sep=1pt},
  shp/.style={font=\scriptsize,text=Muted,inner sep=1pt},
  note/.style={font=\scriptsize,text=Muted,align=center,inner sep=1pt},
  stage/.style={font=\small\bfseries,text=Muted},
  bigop/.style={font=\Large,text=Ink},
  brc/.style={decorate,decoration={brace,amplitude=3pt},draw=Ink!55,line width=.5pt},
}
% \Tens{id}{cx}{cy}{rows}{cols}{color}{depth}{symbol}{shape}
\newcommand{\Tens}[9]{%
  \stepcounter{tseed}%
  \pgfmathsetmacro{\tw}{#5*\U}\pgfmathsetmacro{\th}{#4*\U}%
  \pgfmathsetmacro{\tx}{#2-\tw/2}\pgfmathsetmacro{\ty}{#3-\th/2}%
  \ifnum#7>0
    \foreach \k in {#7,...,1}{%
      \path[draw=#6!45,fill=white,rounded corners=.8pt,line width=.4pt]
        ({\tx+\k*\Off},{\ty+\k*\Off}) rectangle ++(\tw,\th);}%
  \fi
  \foreach \r in {1,...,#4}{\foreach \c in {1,...,#5}{%
    \pgfmathtruncatemacro{\lv}{mod(int(abs(sin(\r*127.1+\c*311.7+\arabic{tseed}*74.7))*1000),7)}%
    \pgfmathtruncatemacro{\tone}{\lv<3?30:(\lv<6?55:80)}%
    \path[fill=#6!\tone,rounded corners=.6pt]
      ({\tx+(\c-1)*\U+.022},{\ty+\th-\r*\U+.022}) rectangle ++({\U-.044},{\U-.044});}}%
  \coordinate (#1-w) at (\tx,#3);
  \coordinate (#1-e) at ({\tx+\tw+#7*\Off},#3);
  \coordinate (#1-n) at (#2,{\ty+\th+#7*\Off});
  \coordinate (#1-s) at (#2,\ty);
  \coordinate (#1-nw) at (\tx,{\ty+\th});
  \coordinate (#1-sw) at (\tx,\ty);
  \node[sym] (#1-sym) at (#2,{\ty-.27}) {$#8$};
  \node[shp] (#1-shp) at (#2,{\ty-.6}) {$#9$};}
% outline one row (r, 1-based) of a tensor whose front face top-left is (x,y)
% \RowMark{x_left}{y_top}{row}{cols}{color}
\newcommand{\RowMark}[5]{%
  \path[draw=#5,line width=1.1pt,rounded corners=1pt]
    ({#1-.03},{#2-#3*\U-.03}) rectangle ++({#4*\U+.06},{\U+.06});}
\newcommand{\ColMark}[5]{% x_left y_top col rows color
  \path[draw=#5,line width=1.1pt,rounded corners=1pt]
    ({#1+(#3-1)*\U-.03},{#2-#4*\U-.03}) rectangle ++({\U+.06},{#4*\U+.06});}
% \TokRow{x_left}{y_center}{roles}: one tile per token, roles in {S,Q,1,2,3,D,G,M,C,P,L,T}
\newcommand{\TokRow}[3]{%
  \foreach \t [count=\i] in {#3}{%
    \path[fill=Tok\t,draw=Ink!25,line width=.3pt,rounded corners=.6pt]
      ({#1+(\i-1)*\U+.022},{#2-\U/2+.022}) rectangle ++({\U-.044},{\U-.044});}}
\newcommand{\TokCol}[3]{% x_center y_top roles (top to bottom)
  \foreach \t [count=\i] in {#3}{%
    \path[fill=Tok\t,draw=Ink!25,line width=.3pt,rounded corners=.6pt]
      ({#1-\U/2+.022},{#2-\i*\U+.022}) rectangle ++({\U-.044},{\U-.044});}}
% \Mask{x_left}{y_top}{n}{groups}{mode}
%  mode 0: causal & (g_j=0 or g_j=g_i)      shared-prefix tree
%  mode 1: all ones                          bidirectional
%  mode 2: causal & (g_j=0 or g_j=g_i or g_i=9)  pointer (9 = decision token)
%  mode 3: causal                            plain causal
%  mode 4: g_j=g_i                           independent bidirectional blocks
%  mode 5: causal with generated tokens (g=5) drawn in Gen colour
\newcommand{\Mask}[5]{%
  \foreach \i in {1,...,#3}{\foreach \j in {1,...,#3}{%
    \pgfmathtruncatemacro{\gi}{{#4}[\i-1]}%
    \pgfmathtruncatemacro{\gj}{{#4}[\j-1]}%
    \ifcase#5
      \pgfmathtruncatemacro{\ok}{and(\j<=\i,or(\gj==0,\gj==\gi))}\or
      \pgfmathtruncatemacro{\ok}{1}\or
      \pgfmathtruncatemacro{\ok}{and(\j<=\i,or(or(\gj==0,\gj==\gi),\gi==9))}\or
      \pgfmathtruncatemacro{\ok}{\j<=\i}\or
      \pgfmathtruncatemacro{\ok}{\gj==\gi}\or
      \pgfmathtruncatemacro{\ok}{\j<=\i}\fi
    \ifnum\ok=1
      \path[fill=Msk!70,rounded corners=.5pt]
        ({#1+(\j-1)*\U+.022},{#2-\i*\U+.022}) rectangle ++({\U-.044},{\U-.044});
    \fi}}%
  \path[draw=Ink!40,line width=.4pt] (#1,#2) rectangle ++({#3*\U},{-#3*\U});}
% integer slot strip: \Slots{x_left}{y_center}{entries}
\newcommand{\Slots}[3]{%
  \foreach \v [count=\i] in {#3}{%
    \node[draw=Idx!70,fill=Idx!12,rounded corners=1pt,minimum width=\U cm-.04cm,
      minimum height=\U cm-.04cm,inner sep=0pt,font=\tiny,text=Ink]
      at ({#1+(\i-.5)*\U},#2) {$\v$};}}
% bottom meaning box: \Meaning{x0}{y_top}{width}{axes}{objects}{mechanism}{id line}
\newcommand{\Meaning}[7]{%
  \node[anchor=north west,fill=Ink!4,rounded corners=3pt,inner sep=9pt,
    text width=#3cm-18pt,align=left,font=\small] (meaning) at (#1,#2) {%
    \begin{tabular}{@{}p{1.1cm}@{\hspace{6pt}}p{\dimexpr #3cm-2.0cm\relax}@{}}
    \textbf{轴} & \raggedright #4\tabularnewline[5pt]
    \textbf{对象} & \raggedright #5\tabularnewline[5pt]
    \textbf{机制} & \raggedright #6\tabularnewline
    \end{tabular}};
  \node[font=\scriptsize,text=Ink!45,anchor=north] at ($(meaning.south)+(0,-.18)$) {#7};}

\begin{document}
\begin{tikzpicture}[x=1cm,y=1cm,text=Ink]
\node[font=\large] at (11.6,4.75)
  {$r_{1:G}\sim\pi_\theta(\cdot\mid x),\quad h_b=H_{b,t^\star_b,:}\ \text{（在 }[x;r]\text{ 上）},\quad
    z_b=W_{\mathrm{lm}}h_b,\quad p_b=\operatorname{softmax}_{k}(z_{b,\ell_k})$};
% ---------- stage A ----------
\node[stage] at (9.0,4.05) {先自回归采样理由（$G$ 次前向），再在末位置按读出 1 读标签分布};
\node[sym,anchor=east] at (1.0,3.05) {$X$};
\node[shp,anchor=east] at (1.0,2.75) {$[B,T]$};
\TokRow{1.2}{2.9}{S,S,S,Q,Q,1,1,2,2,3,3,G,G,G,G,G,D}
\draw[brc] (1.22,3.11) -- (4.48,3.11) node[midway,above=3pt,note] {提示 $x$（$T_0$ 个 token）};
\draw[brc] (4.52,3.11) -- (5.98,3.11) node[midway,above=3pt,note] {理由 $r$（$G$ 步）};
\draw[brc] (6.02,3.11) -- (6.28,3.11) node[midway,above=3pt,note] {$t^\star$};
\Mask{1.2}{2.55}{17}{0,0,0,0,0,0,0,0,0,0,0,0,0,0,0,0,0}{3}
\draw[Gen,line width=.8pt,dashed] (4.5,2.55) -- (4.5,-2.55);
\draw[Gen,line width=.8pt,dashed] (1.2,-0.75) -- (6.3,-0.75);
\node[sym] at (3.75,-2.82) {$M$};
\node[shp] at (3.75,-3.15) {$[T,T]$ 因果；虚线右下为采样阶段};
\node[op,minimum width=2.7cm,minimum height=1.3cm] (dec) at (8.75,0)
  {\textbf{Decoder}\\[2pt]逐 token 采样\\[2pt]KV cache};
\draw[flow] (6.45,0) -- (dec.west);
\draw[flow] (9.6,0.65) .. controls (9.6,1.55) and (7.9,1.55) .. (7.9,0.65);
\node[note] at (8.75,1.6) {$\times G$ 次前向};
\Tens{H}{11.95}{0}{17}{5}{Hid}{1}{H}{[B,T,d]}
\RowMark{11.2}{2.55}{17}{5}{Lgt}
\draw[flow] (dec.east) -- (H-w);
\node[op,minimum width=2.1cm,minimum height=1.1cm] (g1) at (14.75,0)
  {\textbf{Gather}\\[2pt]$H_{b,t^\star_b,:}$};
\draw[flow] (H-e) -- (g1.west);
\Slots{14.45}{-1.25}{16,16}
\node[sym] at (14.75,-1.62) {$t^\star$};
\node[shp] at (14.75,-1.95) {$[B]$};
\draw[dflow] (14.75,-1.08) -- (g1.south);
\Tens{h}{17.4}{0}{2}{5}{Hid}{0}{h}{[B,d]}
\draw[flow] (g1.east) -- (h-w);
% ---------- stage B ----------
\node[stage] at (11.6,-4.1) {读出与读出 1 相同；训练目标覆盖理由 token 或改用强化学习};
\Tens{h2}{1.95}{-6.0}{2}{5}{Hid}{0}{h}{[B,d]}
\node[bigop] at (3.15,-6.0) {$\times$};
\Tens{W}{5.4}{-6.0}{5}{12}{Wgt}{0}{W_{\mathrm{lm}}^{\mathsf T}}{[d,V]}
\Tens{z}{9.75}{-6.0}{2}{12}{Lgt}{0}{z}{[B,V]}
\draw[flow] (W-e) -- (z-w);
\foreach \c in {4,8,11}{\ColMark{7.95}{-5.7}{\c}{2}{Ink!80}}
\node[op,minimum width=1.9cm,minimum height=1.1cm] (g2) at (13.15,-6.0)
  {\textbf{Gather 列}\\[2pt]$z_{b,\ell_k}$};
\draw[flow] (z-e) -- (g2.west);
\Slots{12.7}{-4.75}{3,7,10}
\node[sym,anchor=west] at (13.65,-4.75) {$\ell$ \;{\scriptsize\color{Muted}$[L]$}};
\draw[dflow] (13.15,-4.92) -- (g2.north);
\node[op,minimum width=1.6cm,minimum height=1.1cm] (sm) at (15.4,-6.0)
  {softmax\\[2pt]沿 $L$};
\draw[flow] (g2.east) -- (sm.west);
\Tens{p}{17.4}{-6.0}{2}{3}{Lgt}{0}{p}{[B,L]}
\draw[flow] (sm.east) -- (p-w);
\node[op,minimum width=4.2cm,minimum height=1.3cm] (loss) at (20.95,-6.0)
  {SFT：$[r;y]$ 各位置沿 $V$ 的交叉熵\\[2pt]或 RL：奖励 $R(\hat y,p)$ 的策略梯度\\[2pt]{\scriptsize （如 RLCD / RLCR，奖励含校准项）}};
\draw[flow] (p-e) -- (loss.west);
% ---------- meaning ----------
\Meaning{0.0}{-7.9}{23.4}
 {$B$ 问题数；$T_0$ 提示长度；$G$ 采样的理由长度；$T=T_0+G+1$；$d$ hidden；$V$ 词表；$L$ 标签 token 数。}
 {$r$ 采样得到的理由 token（随机变量，非确定输入）；$t^\star$ 理由之后答案提示的末位置；$\ell$ 标签 token id；$\hat y$ 采样或 argmax 得到的答案；$R$ 奖励。}
 {每问题 $1+G$ 次前向，延迟随 $G$ 线性增长，且结果受采样随机性影响；与 System One“一次前向出判断”的定位冲突，MetisCraft 只作对照或可选兜底。}
 {读出 7：生成推理后读答案分布（对照）@Shichao}
\end{tikzpicture}
\end{document}
